\documentclass[12pt]{article}
\usepackage[a4paper, margin=2cm]{geometry}

\usepackage{fontspec}
\usepackage{polyglossia}
\setdefaultlanguage{russian}
\setmainlanguage{russian}
\setmainfont{Times New Roman}

\begin{document}

\begin{center}
\textbf{\LARGE СЛУЖЕБНАЯ ЗАПИСКА}
\end{center}

\vspace{1cm}

Кому: Ивану Иванову, начальнику отдела по международному экспорту нефти

От кого: Александр Александрович, сотрудник отдела экспорта, г. Санкт-Петербург

Тема: Предложение о корректировке подхода к планированию поставок

Уважаемый Иван Иванов,

В целях повышения точности планирования поставок прошу рассмотреть предложение по переходу от использования усреднённых цен к ориентиру на актуальную рыночную стоимость нефти Brent.

Описание проблемы: текущее использование усреднённых цен приводит к расхождениям между плановыми и фактическими показателями, что увеличивает операционные риски и усложняет бюджетирование.

Предложение: ориентироваться на текущую рыночную стоимость нефти Brent для расчетов планов поставок. Это обеспечит более точное отражение условий рынка и снизит ценовые риски.

Аргумент: на дату 25 марта 2026 года цена Brent составляет примерно 102 USD за баррель. Применение данной котировки позволит корректировать планирование в соответствии с рыночной ситуацией и улучшить точность прогнозов.

Реализация: предлагаю рассмотреть следующие шаги:
- определить частоту обновления котировок (например, ежедневно или еженедельно);
- внедрить процесс пересмотра плана с использованием текущей котировки Brent в каждом квартальном цикле и по мере необходимости;
- обеспечить уведомление заинтересованных сторон о изменениях.

Прошу рассмотреть данный подход и, при отсутствии возражений, утвердить переход к использованию рыночной котировки Brent в процессе планирования.

Готов обсудить детали и предложить конкретный план внедрения.

С уважением,

Александр Александрович

\vfill

\noindent
Дата: 25.03.2026 \hfill Подпись: Александр Александрович

\end{document}